2.26. 下述定理 2.27 在本讲的初版（由 IHES 散发）中曾被表述为一个不正确的猜想。它得益于与 D. S. Rim 的多次讨论，并包含了他的一些结果（D. S. Rim，"Torsion Differentials and Deformations"，定理 2.7(i)）。

%%%%%%%%%% source scan idx 39 | folio 32 %%%%%%%%%%

设 $C$ 是代数闭域 $k$ 上的约化射影曲线，$s\in C(k)$。采用如下记号：

\begin{description}[leftmargin=2.7cm,style=nextline]
\item[$A_s$] $C$ 在 $s$ 处的局部环（或其完备化，这对下文没有影响）；
\item[$\widetilde A_s$] $A_s$ 的正规化；
\item[$\delta(s)$] $\dim_k(\widetilde A_s/A_s)$；
\item[$m_0(s)$] $\operatorname{Spec}(A_s)$ 的形变的完备形式概形 $T$ 的维数。
\end{description}

对 $T$ 的一个不可约分支 $E$，类似地置 $m_0^E(s)=\dim(E)$。

$\widetilde A_s$ 的导子模 $\widetilde\tau$ 及 $A_s$ 的导子模 $\tau$ 均无挠。若 $A_s$ 整，它们嵌入 $A_s$ 的分式域的导子模；一般地，它们嵌入 $A_s$ 的全商环（若干域的积）的导子模。此外，通过这些嵌入，$\tau$ 与 $\widetilde\tau$ 可公度：$\tau/(\tau\cap\widetilde\tau)$ 与 $\widetilde\tau/(\tau\cap\widetilde\tau)$ 在 $k$ 上均为有限维。置

\begin{equation*}
m_1(s)=[\widetilde\tau:\tau]
\mathrel{\mathop=\limits_{\mathrm{def}}}
\dim_k\frac{\widetilde\tau}{\tau\cap\widetilde\tau}
-\dim_k\frac{\tau}{\tau\cap\widetilde\tau}.
\end{equation*}

\underline{定理} 2.27. \underline{设 $E$ 是 $A_s$ 的形变的完备形式概形 $T$ 的一个不可约分支。若 $A_s$ 在 $E$ 的泛点 $e$ 上的形变光滑——也就是说，由 Jacobian 理想定义（VI 5）的、在 $T$ 上有限的形式非光滑子概形的像不包含 $e$——则}

\begin{equation*}
3\delta(s)=m_0^E(s)+m_1(s).
\end{equation*}

%%%%%%%%%% source scan idx 40 | folio 33 %%%%%%%%%%

2.27 的成立与否只依赖于 $A_s$ 的完备化。将 $(C,s)$ 替换为在 $s$ 处与之形式同构的一对 $(C',s)$，可假设下列条件成立：

\begin{enumerate}[label=(\alph*),leftmargin=1.2cm]
\item $s$ 是 $C$ 唯一的奇点；
\item $C$ 没有无穷小自同构。
\end{enumerate}

条件 (b) 保证 $C$ 的形变函子可形式表出。它甚至是有效可形式表出的，因为提升 $C$ 上给定的丰沛可逆层没有障碍。设 $f:X_0\longrightarrow T_0$ 是 $C$ 的万有形变，$t_0$ 是 $T_0$ 的闭点，因而 $C\overset\sim\longrightarrow X_{0t_0}$。

\underline{引理} 2.28. \underline{存在 $k$ 上的有限型概形 $T_1$、点 $t_1\in T_1(k)$ 以及固有平坦态射 $f:X_1\longrightarrow T_1$，其在 $t_1$ 上的纤维为 $C$，使得对每个 $u\in T_1(k)$，$T_1$ 在 $u$ 处的完备化都是纤维 $X_{1u}$ 的形式模概形。}

下面直接证明这个 Artin 一般结果的推论。设 $\mathcal L$ 是 $C$ 上的甚丰沛可逆层，满足 $H^1(C,\mathcal L)=0$，并置 $\mathbb P=\mathbb P(H^0(C,\mathcal L))$。有嵌入 $v:C\hookrightarrow\mathbb P$。设 $T_2$ 是参数化 $\mathbb P$ 中与 $v(C)\subset\mathbb P$ 具有相同 Hilbert 多项式的子簇的 Hilbert 概形，且 $t_2\in T_2$ 对应于 $v(C)$。$T_2$ 在 $t_2$ 处的完备化 $\widehat T_2$ 在 $T_0$ 上光滑。$\widehat T_2/T_0$ 的垂直切丛由 $T_2$ 的切丛的子丛 $L$ 诱导，后者是 $\operatorname{Lie}(\underline{\operatorname{Aut}}(\mathbb P))$ 在 $\underline{\operatorname{Aut}}(\mathbb P)$ 对 $T_2$ 的作用下的像。

在 $t_2$ 的一个邻域中，取方程 $f_1,\ldots,f_n$，其中 $n=\dim\underline{\operatorname{Aut}}(\mathbb P)$，定义 $T_2$ 的一个经过 $t_2$ 的子概形 $T'_1$，且它在 $t_2$ 处的完备化在 $T_0$ 上平展。取 $T_1$ 为 $t_2$ 在 $T'_1$ 中的一个邻域，使其切丛与 $L$ 横截。

%%%%%%%%%% source scan idx 41 | folio 34 %%%%%%%%%%

2.29. 函子态射

\begin{equation*}
(C\text{ 的形变})\longrightarrow(A_s\text{ 的形变})
\end{equation*}

对应于态射 $p:T_0\longrightarrow T$。由 VI 2.10，$p$ 光滑。设 $t$ 是 $T$ 的闭点。$p^{-1}(t)$ 在 $t_0$ 处的 Zariski 切空间分类 $C$ 在 $k[\varepsilon]$ 上的局部平凡形变，其中 $\varepsilon^2=0$。置 $E_0=p^{-1}(E)$，得到

\begin{equation}\tag{2.29.1}
\dim(E_0)=\dim(E)+\dim H^1(C,T^1_{C/k}).
\end{equation}

由于 $X_0/T_0$ 平坦，其所有纤维具有相同的算术亏格 $p$，其中 $1-p=\chi(\underline O)$。根据假设，$E_0$ 的泛点处的纤维光滑。光滑固有曲线的形式模概形光滑且维数为 $3p-3$，故引理 2.28 给出

\begin{equation}\tag{2.29.2}
\dim(E_0)=3p-3.
\end{equation}

设 $\widetilde C$ 是 $C$ 的正规化，其算术亏格为 $\widetilde p=p-\delta(s)$。有

\begin{equation}\tag{2.29.3}
\begin{aligned}
\dim H^1(C,T^1_{C/k})
&=-\chi(C,T^1_{C/k})\\
&=-\chi(\widetilde C,T^1_{\widetilde C/k})+m_1(s)\\
&=[3(p-\delta(s))-3]+m_1(s).
\end{aligned}
\end{equation}

公式（2.29.1）、（2.29.2）和（2.29.3）证明了 2.27。

\underline{猜想} 2.30. \underline{每条约化固有曲线均有光滑形变。}

由 2.27 和 2.28，此猜想蕴含 $A_s$ 的形变的完备形式概形总是等维的，维数为 $3\delta(s)-m_1(s)$。